\subsection{径向伴随与后裔的局部代数}

从正的径向内积选出初级态，需要知道平移和特殊共形生成元怎样互为伴随。共形表示的完整分类可参照《群论》的共形对称与最高权表示章；这里建立后续范数计算所需的关系。

\begin{definition}[径向共形生成元]
取实旋转生成元 $M_{\mu\nu}=-M_{\nu\mu}$、尺度生成元 $D$、平移 $P_\mu$ 和特殊共形生成元 $K_\mu$。它们满足
\begin{align}\label{eq:ghs-conformal-commutators}
 [D,P_\mu]&=P_\mu,&[D,K_\mu]&=-K_\mu,\\
 [K_\mu,P_\nu]&=2\delta_{\mu\nu}D-2M_{\mu\nu},&
 [M_{\mu\nu},P_\rho]&=\delta_{\nu\rho}P_\mu-\delta_{\mu\rho}P_\nu.\notag
\end{align}
在反射正的径向态空间中，球面反演把平移映为特殊共形变换，因此
\begin{equation}\label{eq:ghs-radial-generator-adjoint}
 D^\dagger=D,\qquad P_\mu^\dagger=K_\mu,\qquad M_{\mu\nu}^\dagger=-M_{\mu\nu}.
\end{equation}
$D$ 的本征值无量纲，圆柱物理能量差为 $\hbar c\Delta/R$。$M$ 与带虚数因子的厄密角动量算符采用不同约定。
\end{definition}

\begin{proposition}[幺正界与散度范数]
归一化标量初级态 $|\mathcal O\rangle$ 满足 $K_\mu|\mathcal O\rangle=0$、$D|\mathcal O\rangle=\Delta|\mathcal O\rangle$，则
\begin{equation}\label{eq:ghs-scalar-level-one-gram}
 \langle P_\mu\mathcal O|P_\nu\mathcal O\rangle=2\Delta\delta_{\mu\nu},\qquad
 \|P^2|\mathcal O\rangle\|^2=8d\Delta\left(\Delta-\frac{d-2}{2}\right).
\end{equation}
非恒等标量满足
\begin{equation}\label{eq:ghs-scalar-unitarity-bound}
 \Delta\ge\frac{d-2}{2}.
\end{equation}
对秩 $S\ge1$ 的非零对称无迹初级态，散度后裔满足
\begin{equation}\label{eq:ghs-spin-unitarity-bound-extended}
 \sum\|P^\nu|J_{\nu\mu_2\cdots\mu_S}\rangle\|^2
 =2(\Delta-S-d+2)\sum\|J_{\mu_1\cdots\mu_S}\rangle\|^2,
 \qquad\Delta\ge S+d-2.
\end{equation}
求和按有序张量指标进行；饱和时散度后裔在零范数商空间中消失。
\end{proposition}
\begin{proof}
由 $P^\dagger=K$，一级 Gram 矩阵为 $\langle\mathcal O|[K_\mu,P_\nu]|\mathcal O\rangle$，标量旋转项为零。对迹后裔，逐次交换给出
\begin{align*}
 K_\mu P^2|\mathcal O\rangle
 &=2DP_\mu|\mathcal O\rangle-2M_{\mu\nu}P_\nu|\mathcal O\rangle
   +2P_\mu D|\mathcal O\rangle\\
 &=\bigl[2(\Delta+1)-2(d-1)+2\Delta\bigr]P_\mu|\mathcal O\rangle.
\end{align*}
再左乘 $\langle\mathcal O|K_\mu$ 求和，恢复所述二级范数。

记 $L_{\mu_2\cdots\mu_S}=P^\nu J_{\nu\mu_2\cdots\mu_S}$。旋转生成元在首指标的收缩给出 $(d-1)J_{\alpha\mu_2\cdots}$；其余 $S-1$ 个指标各贡献一个相同张量，含两次重复指标的项由无迹性消去。因此
\[
 K_\alpha|L_{\mu_2\cdots}\rangle=2(\Delta-S-d+2)|J_{\alpha\mu_2\cdots}\rangle.
\]
左乘 $\langle J_{\alpha\mu_2\cdots}|$ 求和得到散度范数恒等式，正性给出下界。
\end{proof}
